% ECONOMICS_PROBLEM: {"id": "D63-3", "title": "Complete EFX allocations for additive goods", "jel_code": "D63", "source_id": "OP-0186"}
% !TeX program = xelatex
\documentclass[11pt,a4paper]{article}
\usepackage[margin=23mm]{geometry}
\usepackage{fontspec}
\setmainfont{Latin Modern Roman}
\usepackage{amsmath,amssymb,mathtools}
\usepackage{xcolor,enumitem,booktabs,longtable,array,xurl,hyperref}
\definecolor{muted}{HTML}{53636C}
\hypersetup{colorlinks=true,urlcolor=blue,linkcolor=blue}
\setlength{\parindent}{0pt}
\setlength{\parskip}{5pt}
\newcommand{\R}{\mathbb R}
\newcommand{\E}{\mathbb E}
\newcommand{\N}{\mathbb N}
\newcommand{\ind}{\mathbf 1}
\newcommand{\Var}{\operatorname{Var}}
\newcommand{\Cov}{\operatorname{Cov}}
\newcommand{\supp}{\operatorname{supp}}
\DeclareMathOperator*{\argmin}{arg\,min}
\DeclareMathOperator*{\argmax}{arg\,max}
\newcommand{\entrymeta}[3]{{\sffamily\small\color{muted}\textbf{\texttt{#1}}\quad|\quad #2\par #3\par}}
\newcommand{\Problemref}[1]{\texttt{#1}}
\newcommand{\JELref}[2]{\texttt{#2} (JEL \texttt{#1})}
\begin{document}
\section*{Complete EFX allocations for additive goods}
\textbf{Problem ID:} \texttt{D63-3}\quad \textbf{JEL:} \texttt{D63}\par
% BEGIN ECONOMICS STATEMENT
\label{problem:OP-0186}
\label{audited:ECON-001}
This entry consolidates three source formulations of complete additive-goods EFX existence. Each formulation and its local sources are retained below.\par
\subsection*{Source formulation: Game theory (GT-056)}

% GAME_THEORY_PROBLEM: {"local_id": "GT-056", "original_number": 56, "kind": "P", "entry_kind": "statement_only", "published": false, "review_required": true, "category": "game_theory"}
\label{game:GT-056}
{\sffamily\small\color{muted}
\textbf{GT-056} | Fair division | \textbf{P}: source-explicit additive-goods existence question.
\par}
\subsection*{Definitions and mathematical statement}
Let $N=[n]$, $n\geq4$, and let $M$ be a finite set of indivisible goods. Agent $i$ has $v_i(S)=\sum_{g\in S}v_{ig}$ with $v_{ig}\geq0$. A complete allocation $A=(A_i)_{i\in N}$ is a partition of $M$, allowing empty bundles. EFX with the positive-good convention requires
\[
 \forall i,j\in N\ \forall g\in A_j,
 \quad v_{ig}>0\ \Longrightarrow\
 v_i(A_i)\geq v_i(A_j\setminus\{g\}).
\]
Determine whether
\[
 \forall n\geq4\ \forall M\text{ finite}\ \forall(v_{ig})\in
 \mathbb R_{\geq0}^{n\times|M|},\quad
 \exists A\text{ complete and EFX}.
\]
\subsection*{Short English statement}
Can all goods always be assigned so that removing any positively valued good eliminates each pairwise envy?
\subsection*{Variants and scope boundaries}
EFX$_0$ checks all removed goods, including zeros, and is a stronger convention. Strict positivity removes the distinction. Partial allocations, approximate EFX, chores, and restricted additive classes are different questions.
\subsection*{Source relationship and status}
The recent source explicitly retains the general additive-goods question. The quantifier includes all $n\geq4$ and makes no assertion that every restricted four-agent family is unresolved. The source's bi-valued results are positive subcases, not a resolution of this general statement.
\subsection*{References}
\begingroup\footnotesize\raggedright
\textbf{[S1]} Zehan Lin, Xiaowei Wu, and Shengwei Zhou (2026 preprint). \emph{When Maximum Nash Welfare Becomes Strongly Fair: Bi-valued Goods}. arXiv:2610.00122v1. \textit{Verified locator:} Section 1.2, ``EFX Allocations,'' retains general additive-goods existence; Section 2 for conventions. \url{https://arxiv.org/html/2610.00122v1}
\par\endgroup

\subsection*{Source formulation: Economic research atlas (EFX supplement)}
\label{atlas:app:efx}
{\sffamily\small\color{muted}\textbf{Supplied example ECON-001} | Supplemental mathematical question\par}
\subsection*{Definitions and assumptions}
Let $n\geq2$, let $N=\{1,\ldots,n\}$, and let $M$ be a finite set of indivisible goods. Agent $i$ has additive valuation $v_i(S)=\sum_{g\in S}v_{ig}$ with $v_{ig}\geq0$. A complete allocation $A=(A_i)_{i\in N}$ is a partition of $M$ into pairwise disjoint, possibly empty bundles. Following the supplied example, call $A$ EFX if
\[
\forall i,j\in N\ \forall g\in A_j:\quad
v_{ig}>0\ \Longrightarrow\ v_i(A_i)\geq v_i(A_j\setminus\{g\}).
\]
This is the convention testing every good positively valued by the \emph{envying} agent. The stronger allocation-level convention EFX$_0$ tests every $g\in A_j$, including those with $v_{ig}=0$. Under strictly positive valuations they coincide.
\subsection*{Formal research question}
Determine whether the following universal existence assertion is true:
\[
\forall n\geq4\ \forall M\text{ finite}\ \forall (v_{ig})\in\mathbb R_{\geq0}^{n\times|M|},
\quad \exists A\text{ a complete EFX allocation}.
\]
A negative answer must exhibit a finite additive-goods instance with no complete EFX allocation; a positive answer must cover every instance in the stated domain. Polynomial-time computation is a separate requirement and is not included in this existence assertion.
\subsection*{Interpretation and scope}
The question asks whether all goods can be assigned while every envy comparison disappears after removing any positively valued good from the envied bundle. The focus on $n\geq4$ removes known small-agent benchmarks reported by the status source. Approximate EFX, partial allocations and nonadditive valuations change the target. Chore costs require a different inequality and removal from the envious agent's own bundle, so chore results cannot be transferred directly.
\subsection*{Source audit and status}
Both supplied references are identifiable. The version-pinned June 2026 preprint explicitly retains additive-goods EFX as open in Section 6 and distinguishes its chore results. This is a source-date report, not an independent proof review or exhaustive certification through 8 October 2026. The supplied definition is kept explicit to avoid ambiguity about zero-valued goods. The supplement is additional to the 130 atlas entries and is not given an atlas identifier.
\subsection*{References}
{\footnotesize\raggedright
\begin{enumerate}[label={\textbf{[S\arabic*]}},leftmargin=3em]
\item Wentao He and Biaoshuai Tao (2026). \emph{EFX for Additive Chores: Nonexistence, Pareto Incompatibility, and Bi-Valued Existence}. arXiv:2606.08872v1, 7 June 2026, preprint.\par
\textit{Locator and source role:} Introduction and Section 6, additive-goods open-status discussion; source-date statement only.\par
\url{https://arxiv.org/html/2606.08872v1}
\item Georgios Amanatidis, Haris Aziz, Georgios Birmpas, Aris Filos-Ratsikas, Bo Li, Hervé Moulin, Alexandros A. Voudouris, and Xiaowei Wu (2023). \emph{Fair division of indivisible goods: Recent progress and open questions}. Artificial Intelligence 322, 103965. DOI 10.1016/j.artint.2023.103965.\par
\textit{Locator and source role:} Survey's EFX discussion; verifies bibliographic identity and the established research context.\par
\url{https://arxiv.org/abs/2208.08782}
\end{enumerate}}

\subsection*{Source formulation: Formal economic research definitions (30007558)}
\label{definitions:problem:30007558}
\textbf{Audit classification:} Explicit published open question. The additive complete-allocation question is explicitly posed in 2016 and remains stated as open in 2026 primary sources. Broader monotone/submodular existence is now negatively resolved in 2026 preprints.
\subsection*{Definitions and precise research target}
\textbf{Model and research target.} For every positive integer \(n\), let \(N=\{1,\ldots,n\}\) be agents and let \(G\) be any finite set of indivisible goods. Agent \(i\) has an additive nonnegative valuation \(v_i\), specified by numbers \(v_i(g)\geq0\) for \(g\in G\), with \(v_i(S)=\sum_{g\in S}v_i(g)\) for every bundle \(S\subseteq G\). A complete allocation is a tuple \(A=(A_1,\ldots,A_n)\) of pairwise disjoint bundles whose union equals \(G\); empty bundles are allowed. Under the original positive-good convention, require
\[
 \forall i,j\in N\ \forall g\in A_j\text{ with }v_i(g)>0:
 \quad v_i(A_i)\geq v_i(A_j\setminus\{g\}).
\]
Does a complete allocation satisfying these inequalities exist for every such instance? Resolve the universal existence question by proving the assertion for arbitrary \(n\), \(G\), and valuations, or by exhibiting a finite counterexample with a proof that every complete allocation violates an inequality. No computational efficiency requirement is imposed on an existence proof. Zero-valued goods require care: the displayed condition tests only goods positively valued by the potentially envious agent. This is frequently called \(\mathrm{EFX}^{+}\); some later authors reserve EFX or \(\mathrm{EFX}^{0}\) for the stronger condition testing all goods, including zero-valued ones.

\emph{Definitions and maintained assumptions (editorial unless a source definition is named).} For \(v_i(g)\ge0\), define \(v_i(S)=\sum_{g\in S}v_i(g)\) and \(v_i(\varnothing)=0\). Complete allocations partition \(G\), with empty bundles allowed. \(\mathrm{EFX}^{+}\) tests only \(g\in A_j\) with \(v_i(g)>0\); \(\mathrm{EFX}^{0}\) tests every \(g\in A_j\). Each uses weak inequality \(v_i(A_i)\ge v_i(A_j\setminus\{g\})\). For a nonempty finite bundle the latter is equivalent to testing its least-valued good. An empty comparison bundle has no test. Both conventions coincide for strictly positive item values. Existence is a universal quantifier over all finite sizes and all real nonnegative valuation matrices; exhaustive search decides each rational instance but does not prove universal existence. Three additive agents are a solved case. The August 2026 Alkassar--Fouz--Mehlhorn preprint reports a computer-assisted existence result for four agents and at most nine goods; this audit verifies the reported statement, not its full certificate corpus.
\subsection*{Variants and assumption changes}
\begin{enumerate}
\item \textbf{Zero-tolerant additive version (same open family).} Replace the positive-good test by \(\mathrm{EFX}^{0}\). To transfer an existence result from strictly positive values, add \(\epsilon>0\) to every item value and extract a constant-allocation subsequence as \(\epsilon\downarrow0\); finiteness of allocations and weak inequalities justify the transfer.
\item \textbf{Nonadditive extension (resolved negative, 2026 preprint).} Allow normalized monotone submodular valuations, meaning \(v_i(\varnothing)=0\), monotonicity, and \(v_i(S)+v_i(T)\ge v_i(S\cup T)+v_i(S\cap T)\). Mackenzie--Suzuki Theorem 1.1 gives a three-agent eight-good counterexample under the all-good convention. This extension must not be listed as an open universal-existence conjecture.
\end{enumerate}
\subsection*{References and source audit}
\begin{enumerate}
\item Definition4.4 and following paragraph, p.317: original additive complete-allocation question. \url{https://eprints.gla.ac.uk/123283/1/123283.pdf}
\item Abstract: general additive question still stated open; restricted hypergraphs solved. \url{https://ojs.aaai.org/index.php/AAAI/article/view/38759}
\item Abstract/main theorem: three additive agents solved. \url{https://arxiv.org/pdf/2002.05119}
\item Abstract/§1: four-agent \ensuremath{\leq}9-good result reported in August2026 preprint; global additive question remains. \url{https://arxiv.org/pdf/2608.08590}
\item Theorem1.1: submodular EFX0 universal existence negatively resolved in May2026 preprint. \url{https://arxiv.org/pdf/2605.06451}
\end{enumerate}
\begin{enumerate}\raggedright
\item\hypertarget{definitions:source:30007558:1}{} Ioannis Caragiannis; David Kurokawa; Hervé Moulin; Ariel D. Procaccia; Nisarg Shah; Junxing Wang. 2016. \emph{The Unreasonable Fairness of Maximum Nash Welfare}. §4.2, Definition 4.4 and following paragraph, proceedings p.317.
\url{https://eprints.gla.ac.uk/123283/1/123283.pdf}.
DOI: \url{https://doi.org/10.1145/2940716.2940726}.
\item\hypertarget{definitions:source:30007558:2}{} Thanasis Lianeas; Alkmini Sgouritsa; Minas Marios Sotiriou. 2026. \emph{EFX Allocation in (Multi)Hypergraphs}. Abstract.
\url{https://ojs.aaai.org/index.php/AAAI/article/view/38759}.
DOI: \url{https://doi.org/10.1609/aaai.v40i20.38759}.
\item\hypertarget{definitions:source:30007558:3}{} Bhaskar Ray Chaudhury; Jugal Garg; Kurt Mehlhorn. 2020. \emph{EFX Exists for Three Agents}. Abstract; main existence theorem.
\url{https://arxiv.org/pdf/2002.05119}.
\item\hypertarget{definitions:source:30007558:4}{} Eyad Alkassar; Mahmoud Fouz; Kurt Mehlhorn. 2026. \emph{Complete EFX Allocations Exist for Four Additive Agents and Up to Nine Goods}. Abstract; §1 and main theorem, arXiv v1 August9,2026.
\url{https://arxiv.org/pdf/2608.08590}.
\item\hypertarget{definitions:source:30007558:5}{} Simon Mackenzie; Mashbat Suzuki. 2026. \emph{Counterexamples to EFX for Submodular and Subadditive Valuations}. Theorem1.1; section3.2 and Theorem3.8 (submodular counterexample); Theorem1.2 (approximate subadditive).
\url{https://arxiv.org/pdf/2605.06451}.
\end{enumerate}

\subsection*{Common conventions: Source mathematical conventions}

\textbf{Finite games.} For a finite set $A$, $\Delta(A)$ is its probability
simplex. Players choose independent mixed strategies in Nash equilibrium;
correlation is allowed only when explicitly part of the solution concept.
In a game with payoff functions $u_i$ and strategy profile $x$, an additive
$\varepsilon$ Nash equilibrium satisfies
\[
 u_i(x)\geq u_i(x_i',x_{-i})-\varepsilon
 \quad\text{for every player $i$ and every admissible deviation $x_i'$}.
\]
For numerical approximation, payoffs are normalized as locally specified.
An $\varepsilon$ payoff residual is distinct from distance at most
$\varepsilon$ to the set of exact equilibria. Point convergence, convergence
of empirical play, and convergence in distance to an equilibrium set are
also distinct.

\textbf{Computational representations.} Unless stated otherwise, a complexity
question takes rational data with binary encoding; polynomial time is measured
in total bit length. A fixed-accuracy polynomial algorithm is not necessarily
polynomial in $1/\varepsilon$, and neither is necessarily polynomial in
$\log(1/\varepsilon)$. Succinct games and value, demand, or sampling oracles
must declare their interfaces. Exact real solutions require an explicit
representation; an unspecified list of arbitrary real numbers is not a
finite computational output. A claimed algorithmic barrier is meaningful
only for its specified model and reductions.

\textbf{Dynamic games.} Histories, information, observation, and allowable
randomization are part of the model. A stationary strategy depends only on
the current state; a positional strategy in a deterministic graph game
chooses one action at each controlled vertex. Nash equilibrium at an initial
state and equilibrium after every history are different requirements.
An expectation of a pathwise $\liminf$ payoff, a $\liminf$ of expected averages,
a limiting expected average, and a uniform finite-horizon equilibrium are
different criteria. None is silently substituted for another.

\textbf{Allocation and mechanisms.} A complete allocation assigns every item
exactly once unless otherwise stated. Zero-valued goods, negative values,
capacity constraints, feasibility, lotteries, and copies of goods affect
fairness definitions and are specified locally. Dominant-strategy incentive
compatibility, Bayesian incentive compatibility, universal truthfulness,
and truthfulness in expectation impose different inequalities. Whenever
an objective ratio has zero denominator, its convention is stated or those
instances are excluded.

\textbf{Infinite-dimensional models.} Expectations, integrals, stochastic
equations, and optimization objectives require their local regularity,
integrability, filtrations, and admissible domains. A backward--forward
mean-field system is a boundary-value problem; it is not an initial-value
semiflow without an additional construction. Long-time, large-population,
and vanishing-noise limits require an order of limits and a convergence
topology. If a minimum need not be attained, the objective is its infimum
or supremum and the approximation requirement is stated separately.

\subsection*{Common conventions: Additional-source status and mathematical conventions}

\label{audited:conventions}

Of the 100 supplied ECON entries, 65 distinct problems appear here; 32 close
matches refine earlier entries, and three duplicate questions map to existing
entries. Appendix~\ref{app:audited-crosswalk} lists every destination.
The attachment's P/R/S/U distinctions and source audits are retained. Its
precise mathematical formalizations are editorial unless the individual audit
says otherwise. Candidate subsidy targets ECON-004--006 retain their U flags;
the resolved approval-core question ECON-007 updates OP-0202 above.
The following supplied conventions govern both this part and the attributed
ECON refinements elsewhere in the catalogue, subject to local restrictions.

\textbf{Parameterized questions.} A problem family lists inputs that must be fixed before an instance is evaluated: population, horizon, policies, information, data law, model class and welfare weights. These are required choices, not quantities the citations are assumed to specify. Undefined applications should not be treated as identified estimands.

\textbf{Indices and integrability.} Sets of agents, firms and regions are finite unless a population measure is explicitly used. \(i,j\) index units, \(r\) regions, \(t\) dates and \(h\) horizons. \(\mathbb E\) and \(\Pr\) refer to the declared population and law; \(\mathbf1\{A\}\) is an indicator. Every displayed expectation and discounted sum needs existence in the stated domain. Infinite horizons use a discount in \((0,1)\) and a convergence condition; finite horizons may use discount one.

\textbf{Optimization and existence.} A displayed \(\max\), \(\min\), or \(\arg\max\) is conditional on attainment. For finite feasible menus this is immediate; general spaces need continuity and compactness/coercivity or another existence argument. Without attainment, the target is the supremum/infimum and arbitrarily near-optimal feasible choices. \(\inf\varnothing=+\infty\). Empty feasibility and an unidentified target are different issues.

\textbf{Potential outcomes.} \(Y_i(z)\) is the outcome under a specified intervention, not the observed conditional mean among people choosing \(z\). In binary comparisons \(\Delta Y_i=Y_i(1)-Y_i(0)\). Specify assignment, treatment versions, compliance, information and observation horizon. Interference is allowed under a full policy vector; compressed exposure notation needs a justified mapping. No interference, consistency and overlap are substantive assumptions, not automatic consequences of notation.

\textbf{Populations and observation.} Fix baseline eligibility and population weights. Conditioning on post-policy employment, survival, relocation or program uptake can change the population being compared. Death is not ordinary loss to follow-up; outcomes truncated by death need a composite convention, joint survival target, or explicitly defined survivor stratum. Fertility-changing interventions need a population functional rather than an unexplained fixed descendant cohort. Measurement and missingness models must be supplied.

\textbf{Identification.} A structural model \(m\in\mathcal M\) implies observable law \(P_m\) and target \(\theta(m)\). Identification means
\[
 P_{m_1}=P_{m_2}\ \Longrightarrow\ \theta(m_1)=\theta(m_2)
 \quad\text{for all }m_1,m_2\in\mathcal M .
\]
Without identification, the target is
\[
 \Theta(P;\mathcal M)=\{\theta(m):m\in\mathcal M,\ P_m=P\}.
\]
Writing \(\theta(P)\) before establishing identification can hide incompatible latent models. Confidence sets additionally need a sampling law, confidence level, asymptotic sequence and uniformity class. Point identification, coverage and informative precision are separate properties.

\textbf{Mediation.} Total effects of randomized assignment are distinct from cross-world natural indirect effects. Belief/effort-channel effects require a meaningful mediator intervention and additional measurement, exclusion or mediation assumptions. Randomization of the initial policy alone does not supply those assumptions. Report bounds or interventional alternatives when the stronger target is unsupported.

\textbf{Equilibrium.} A counterfactual \(W(z)\), \(p(z)\), or \(\mu_t^z\) requires individual optimization, feasibility, government/resource budgets, market clearing and state-transition laws. State equilibrium existence and selection, or report ranges over compatible equilibria. Initial-state integration includes subsequent endogenous transitions; current-state integration must use the policy-dependent distribution. Reduced-form invariance after policy is not assumed.

\textbf{Welfare accounts.} Scores, utility, health, dollars and ecological quantities require explicit conversion before addition. Fees, taxes, subsidies, wages and extortion payments are transfers between parties; distributional weights can make them welfare relevant, but their face value is not automatically a resource loss. State ownership, geographic boundaries, foreign-income weights and fiscal finance. Do not add profits to household welfare already including dividends, or count land capitalization and the underlying benefit twice. A benefit-cost expression must distinguish gross benefits from net welfare and subtract every real cost exactly once.

\textbf{Derivatives and risk.} Log derivatives require positive arguments; local responses require admissible perturbations and specified equilibrium branches. Differentiation under expectations needs regularity/domination, and boundaries may allow only one-sided derivatives. Risk uses a specified law; ambiguity uses a declared nonempty law set on a common history space. Adaptive policies must be nonanticipative. A robust criterion is an editorial choice unless attributed explicitly.

\textbf{Scope overlaps.} Allocation, collective choice, contracts, household bargaining and coordinated policy overlap economics with optimization and game theory. They are retained because they appeared in the original catalogue; no claim of a strictly game-theory-free selection is made.
% END ECONOMICS STATEMENT
\end{document}
